\documentclass[11pt,a4paper]{article}
\usepackage[utf8]{inputenc}
\usepackage[T1]{fontenc}
\usepackage{amsmath,amssymb,amsthm}
\usepackage{listings}
\usepackage{geometry}
\usepackage{hyperref}
\geometry{margin=2.5cm}

\lstset{
  language=C++,
  basicstyle=\ttfamily\small,
  keywordstyle=\bfseries,
  breaklines=true,
  numbers=left,
  numberstyle=\tiny,
  frame=single,
  tabsize=2
}

\title{IOI 2021 -- Registers}
\author{Solution}
\date{}

\begin{document}
\maketitle

\section{Problem Summary}
You have $m = 100$ registers, each containing $b = 2000$ bits. Register 0 initially contains $n$ values, each $k$ bits wide, packed consecutively. The remaining bits are 0. You must output the sorted values (or just the minimum, depending on the subtask) in register 0 using a limited number of instructions:

Instructions: AND, OR, XOR, NOT (bitwise on full registers), LEFT shift, RIGHT shift, ADD (modular addition of two registers).

\section{Solution Approach}

\subsection{Subtask 1: Finding Minimum}
For finding the minimum of $n$ values of $k$ bits each:

\subsubsection{Pairwise Comparison}
To compare two $k$-bit numbers $A$ and $B$: compute $A - B$ (using ADD with $B$'s complement). The sign bit (bit $k-1$ of the result) tells us which is larger.

\subsubsection{Tournament Approach}
\begin{enumerate}
  \item Copy register 0. Shift the copy right by $k$ bits to align adjacent values.
  \item Compare pairs and keep the minimum.
  \item Repeat $\lceil \log_2 n \rceil$ times.
\end{enumerate}

\subsection{Subtask 2: Full Sort (Sorting Network)}
Implement a sorting network (e.g., odd-even merge sort or bitonic sort) using register operations.

Each comparator (compare-and-swap of two values):
\begin{enumerate}
  \item Extract the two values into separate registers.
  \item Compare them (subtract and check sign).
  \item Based on the comparison, write $\min$ and $\max$ back.
\end{enumerate}

A bitonic sort network of $n$ elements uses $O(n \log^2 n)$ comparators. Each comparator is $O(1)$ register operations.

\section{C++ Solution}

\begin{lstlisting}
#include <bits/stdc++.h>
using namespace std;

// Grader-provided functions (declarations)
void append_move(int t, int x);       // r[t] = r[x]
void append_store(int t, vector<bool> v); // r[t] = constant v
void append_and(int t, int x, int y); // r[t] = r[x] & r[y]
void append_or(int t, int x, int y);  // r[t] = r[x] | r[y]
void append_xor(int t, int x, int y); // r[t] = r[x] ^ r[y]
void append_not(int t, int x);        // r[t] = ~r[x]
void append_left(int t, int x, int p);  // r[t] = r[x] << p
void append_right(int t, int x, int p); // r[t] = r[x] >> p
void append_add(int t, int x, int y); // r[t] = (r[x] + r[y]) mod 2^b

const int B = 2000;

void construct_instructions(int s, int n, int k, int q) {
    if (s == 0) {
        // Find minimum
        // Strategy: pairwise tournament

        // Create mask: k bits of 1s
        vector<bool> mask_k(B, false);
        for (int i = 0; i < k; i++) mask_k[i] = true;

        // Store one in register for subtraction
        vector<bool> one(B, false);
        one[0] = true;

        append_store(2, mask_k); // mask for k bits

        for (int step = 1; step < n; step *= 2) {
            // r[0] has values at positions 0, step, 2*step, ...
            // Shift right by step*k to align pairs
            append_right(1, 0, step * k); // r[1] = r[0] >> (step*k)

            // Compare r[0] and r[1] at each k-bit slot
            // For each slot: if r[0] > r[1], keep r[1]; else keep r[0]

            // Compute r[0] - r[1] to find comparison
            // We need to do this per-slot (k-bit comparison)
            // Simple approach: for each pair of values, compute min

            // Extract first value from each pair
            // Mask both to k bits (per slot)
            // Actually, we need per-slot min. This is complex for parallel.

            // For the simple subtask: just iterate and take pairwise min

            // r[3] = r[0] AND mask (extract first value)
            append_and(3, 0, 2);
            // r[4] = r[1] AND mask (extract second value)
            append_and(4, 1, 2);

            // Compare: r[5] = r[3] - r[4] = r[3] + NOT(r[4]) + 1
            append_not(5, 4);
            append_store(6, one);
            append_add(5, 5, 6);
            append_add(5, 3, 5);

            // Check sign bit (bit k-1 of each slot): if 1, r[3] < r[4] (underflow)
            // Actually, if r[3] < r[4], then r[3] - r[4] wraps around and bit (k) would be 0
            // For k-bit values: r[3] - r[4] mod 2^k. If r[3] >= r[4], result = r[3]-r[4].
            // If r[3] < r[4], result = 2^k - (r[4]-r[3]), which has bit k-1 set (if k >= 2).
            // So bit k-1 of (r[3] - r[4]) mod 2^b:
            // We need to check bit k of the add result before mod 2^k...
            // Actually, let's use a different approach.

            // Subtract and check: compute r[4] - r[3]. If result's MSB (bit k-1) is set,
            // then r[4] < r[3] (with unsigned wraparound interpretation)... no.

            // For unsigned comparison: r[3] > r[4] iff r[3] - r[4] (mod 2^b) in [1, 2^(b-1)-1]
            // But we only care about k-bit values.

            // Better: compute r[3] + (2^k - r[4]) = r[3] - r[4] + 2^k.
            // If r[3] >= r[4]: result >= 2^k, so bit k is set.
            // If r[3] < r[4]: result < 2^k, so bit k is clear.

            // Compute 2^k - r[4] = NOT(r[4]) + 1 (k-bit complement), but we already have this.
            // r[5] = r[3] + (2^k - r[4]) mod 2^b
            // Actually: append_not(5, 4) gives ~r[4] in all B bits.
            // We want only the k-bit complement: mask first.
            append_and(4, 1, 2); // r[4] = second value, k bits
            append_not(5, 4);
            append_and(5, 5, 2); // k-bit NOT
            append_add(5, 5, 6); // 2^k - r[4] (k-bit complement), but this gives 2^k if r[4]=0
            append_add(5, 3, 5); // r[3] + 2^k - r[4]

            // Bit k of r[5]: 1 if r[3] >= r[4], 0 if r[3] < r[4]
            // If bit k = 1: min = r[4]. If bit k = 0: min = r[3].
            append_right(7, 5, k); // shift bit k to bit 0
            append_and(7, 7, 6);   // isolate bit 0
            // r[7] = 1 if r[3] >= r[4], else 0

            // Spread to k-bit mask: subtract 1, then NOT, then AND with mask
            // If r[7] = 1: r[7]-1 = 0, NOT = all 1s -> mask = all 1s in k bits
            // If r[7] = 0: r[7]-1 = -1 = all 1s, NOT = 0 -> mask = 0 in k bits

            // Build selection mask
            // r[7] = 0 or 1. We want a k-bit mask that's all-1s if r[3] < r[4].
            // r[3] < r[4] iff r[7] = 0.
            // Negate r[7]: if 0 -> 1, if 1 -> 0
            append_xor(7, 7, 6); // flip bit 0
            append_and(7, 7, 6); // r[7] = 1 if r[3] < r[4], else 0
            // Spread: r[7] - 1 = 0 or -1
            // Actually: multiply by mask. Use: left shift and OR repeatedly.
            // Or: subtract, then AND with mask.

            // Simpler: use NOT and ADD to create mask
            // r[7] = 0 or 1. subtract 1: if 1 -> 0, if 0 -> all 1s.
            append_not(8, 6); // all 1s except bit 0 = -2 in two's complement
            append_add(8, 7, 8); // r[7] - 1 = r[7] + (-1) ... no.
            // r[7] + NOT(1) + 1... this is getting complex.

            // Even simpler: if r[3] < r[4], pick r[3]; else pick r[4].
            // r[0] = min(r[3], r[4]) = r[3] if r[3]<r[4] else r[4]
            // Use conditional: min = r[3] XOR ((r[3] XOR r[4]) AND mask)
            // where mask = all 1s if we should pick r[4].

            // Let's just use: negate r[7] (0 or 1 -> 0...0 or 1...1 k-bit mask)
            // r[7] = 1 iff r[3] < r[4] (select r[3])
            // We want: r[0] = r[3] if r[7]=1, r[4] if r[7]=0
            // = (r[3] AND expanded(r[7])) OR (r[4] AND expanded(NOT r[7]))

            // Expand r[7] to k bits: left shift and OR
            append_move(8, 7);
            for (int bit = 0; (1 << bit) < k; bit++) {
                append_left(9, 8, 1 << bit);
                append_or(8, 8, 9);
            }
            append_and(8, 8, 2); // r[8] = k-bit mask

            append_and(9, 3, 8);  // r[3] if selected
            append_not(10, 8);
            append_and(10, 10, 2);
            append_and(10, 4, 10); // r[4] if not selected
            append_or(0, 9, 10);   // r[0] = min(r[3], r[4])
        }
    } else {
        // Full sort: implement bitonic sort or odd-even sort
        // For simplicity, implement bubble sort with O(n^2) comparators

        for (int i = 0; i < n; i++) {
            for (int j = 0; j < n - 1 - i; j++) {
                // Compare-and-swap positions j and j+1
                // Extract value at position j (bits j*k to (j+1)*k-1) and j+1

                vector<bool> mask_j(B, false), mask_j1(B, false);
                for (int b = 0; b < k; b++) {
                    mask_j[j * k + b] = true;
                    mask_j1[(j + 1) * k + b] = true;
                }
                append_store(2, mask_j);
                append_store(3, mask_j1);

                append_and(4, 0, 2);           // value at j (in position j)
                append_right(4, 4, j * k);     // shift to position 0
                append_and(5, 0, 3);           // value at j+1 (in position j+1)
                append_right(5, 5, (j+1) * k); // shift to position 0

                // Compare r[4] and r[5] (both at position 0, k bits)
                vector<bool> one_v(B, false);
                one_v[0] = true;
                append_store(6, one_v);

                vector<bool> kmask(B, false);
                for (int b = 0; b < k; b++) kmask[b] = true;
                append_store(7, kmask);

                // r[4] - r[5] + 2^k: check bit k for comparison
                append_not(8, 5);
                append_and(8, 8, 7);
                append_add(8, 8, 6); // 2^k - r[5]
                append_add(8, 4, 8); // r[4] - r[5] + 2^k

                // Bit k: 1 if r[4] >= r[5] (need swap), 0 if r[4] < r[5] (no swap)
                append_right(9, 8, k);
                append_and(9, 9, 6); // r[9] = 1 if need swap

                // If swap: r[0] = r[0] with positions j and j+1 swapped
                // min goes to position j, max goes to position j+1

                // Build XOR mask for swap
                append_xor(10, 4, 5); // diff = r[4] XOR r[5]

                // Expand r[9] to k bits
                append_move(11, 9);
                for (int bit = 0; (1 << bit) < k; bit++) {
                    append_left(12, 11, 1 << bit);
                    append_or(11, 11, 12);
                }
                append_and(11, 11, 7); // k-bit swap mask

                append_and(10, 10, 11); // conditional diff

                // XOR into both positions
                append_left(12, 10, j * k);
                append_left(13, 10, (j + 1) * k);
                append_xor(0, 0, 12);
                append_xor(0, 0, 13);
            }
        }
    }
}
\end{lstlisting}

\section{Complexity Analysis}
\begin{itemize}
  \item \textbf{Subtask 1 (Minimum):} $O(k \log n)$ instructions for the tournament approach with $\log n$ rounds.
  \item \textbf{Subtask 2 (Full Sort):} $O(n^2 k)$ instructions for bubble sort. An optimized bitonic sort network would use $O(n \log^2 n \cdot k)$ instructions.
  \item \textbf{Register count:} Uses approximately 15 registers.
\end{itemize}

\end{document}
